\documentclass{situationist-preprint}

\title{There Is No Pure Datum}
\msid{PHE-02-b}
\msfamily{PHE-02}
\msdiscipline{Phenomenology}
\mstype{Consequence}
\mskeywords{givenness; sense data; myth of the given; hyle; pure givenness}

\begin{document}
\maketitle

\begin{abstract}
PHE-02-a argues that givenness is an irreducibly three-place relation: something is given to someone as something. We derive the consequence that there is no pure datum, that is, nothing given to a subject without being given as anything. The sense-datum, the hyletic content of Husserl's analyses taken as given, and Marion's pure givenness are each either given under some sense, however thin, or not given at all. We distinguish this conclusion from Sellars's attack on the myth of the given, which it supports by a different route, and answer the objection that a sufficiently thin sense is as good as none.
\end{abstract}

\section{The derivation}

Call an item $x$ a \emph{pure datum} for a subject $s$ if $x$ is given to $s$ and there is no sense $F$ such that $x$ is given to $s$ as $F$. In the notation of Definition~2.1 of PHE-02-a, $x$ is a pure datum for $s$ if $x$ is given to $s$ while $G(x,s,F)$ fails for every $F$.

\begin{proposition}\label{p:nopure}
There is no pure datum.
\end{proposition}

\begin{proof}
Givenness is the relation $G$, and by Theorem~3.4 of PHE-02-a it is irreducibly three-place: there is no two-place relation of item and subject on which its truth depends alone (Lemma~3.1 there, with Definition~2.2). Hence \enquote{$x$ is given to $s$} has no content except as the existential generalization \enquote{for some $F$, $G(x,s,F)$}. If $x$ is given to $s$, some such $F$ exists; if none exists, $x$ is not given to $s$.
\end{proof}

The proof is short because the work is done in PHE-02-a. What the consequence adds is a classification of the candidates. A purported pure datum either (i) is given under a sense that the theory fails to mention, or (ii) is not given but stands in some other relation to the subject, such as causing her states or being lived through by her.

The sense-datum of early analytic philosophy belongs to case (i): a red patch is given as red, as extended, as there, and the theory's own descriptions supply the sense it officially omits. Husserl's hyletic data belong to case (ii): in \emph{Ideas I} the sensuous contents are not objects of perception but are lived through and animated by apprehension, which gives the thing perceived through them \citep[\S85]{Husserl1913}. Taken as data in their own right, they are no longer hyle but objects of a reflective act, and are then given as sensations. Marion's pure givenness, which precedes every determination \citep{Marion1997}, belongs to case (ii) on the reading discussed in PHE-02-a: what precedes every \enquote{as} is not given but is the condition of being given, and to call it given is to change the subject.

\section{Objections}

\begin{objection}
A sense such as \enquote{something there} is so thin that the datum given under it is pure in every sense that matters.
\end{objection}

\begin{reply}
Thin is not none. \enquote{Something there} locates, and the same item can fail to be given as there, for instance when it is given in memory as having been there. The minimal pairs of PHE-02-a apply at every degree of thinness: the same item is given to the same subject under a thin sense and under a richer one, and the difference is a difference in givenness. A theory of the datum that needed only thin senses would still need senses.
\end{reply}

\begin{objection}
Sellars has already shown that the given is a myth \citep{Sellars1956}. The present result adds nothing.
\end{objection}

\begin{reply}
The conclusions overlap, but the routes differ. Sellars argues that nothing can be both non-inferentially known and independent of conceptual capacities, so that the given cannot serve as a foundation of knowledge. The present argument concerns givenness itself, not its epistemic role, and does not assume that every sense is conceptual: an aspect of perception, such as a profile seen from the left, is a sense in the required meaning. Hence the conclusion holds even for those who reject Sellars's conceptualism.
\end{reply}

\section{Conclusion}

There is givenness, but there is no pure datum. Whatever is given is given as something, and what is not given as anything is not given. The search for a pure datum is therefore a search for something other than givenness, and the theories that claim to find it have found causes, conditions, or lived contents.

\begin{thebibliography}{99}\small
\bibitem[Husserl(1913)]{Husserl1913} Husserl, E. (1913). Ideen zu einer reinen Phänomenologie und phänomenologischen Philosophie. Erstes Buch. \emph{Jahrbuch für Philosophie und phänomenologische Forschung} 1.
\bibitem[Marion(1997)]{Marion1997} Marion, J.-L. (1997). \emph{Étant donné. Essai d'une phénoménologie de la donation}. Presses Universitaires de France.
\bibitem[Sellars(1956)]{Sellars1956} Sellars, W. (1956). Empiricism and the philosophy of mind. In H. Feigl and M. Scriven (eds.), \emph{Minnesota Studies in the Philosophy of Science}, vol.~1. University of Minnesota Press.
\end{thebibliography}

\end{document}
